% TOP_PROBLEM: {"id": 3245, "title": "List chromatic number and maximum degree of bipartite graphs", "category_id": 3, "category": {"id": 3, "name": "graph_theory", "display_name": "Graph Theory", "description": "Problems involving graphs, networks, and their properties.", "slug": "graph-theory", "order_index": 3, "created_at": "2026-07-31T15:26:25.671Z"}, "status": "open", "problem_number": "OPG-47358", "set_id": 12, "statusQualification": "The logarithmic statement is interpreted for maximum degree at least two; an equivalent all-degree formulation is supplied."}
% FIRST_POSTED: 2026-10-01
% ENTRY_KIND: statement_only
% UnsolvedMath ID 3245; source catalogue code OPG-47358
% !TeX program = lualatex
% Definition and sources only; this file is not an LLM solution attempt.
\documentclass[11pt,a4paper]{article}
\usepackage[margin=25mm]{geometry}
\usepackage{fontspec}
\setmainfont{lmroman10-regular.otf}[BoldFont=lmroman10-bold.otf,
 ItalicFont=lmroman10-italic.otf,BoldItalicFont=lmroman10-bolditalic.otf]
\usepackage{amsmath,amssymb,mathtools}
\usepackage{unicode-math}
\setmathfont{latinmodern-math.otf}
\AtBeginDocument{\let\setminus\smallsetminus}
\usepackage{xurl,hyperref}
\hypersetup{colorlinks=true,urlcolor=blue}
\setcounter{secnumdepth}{0}
\setlength{\parindent}{0pt}
\setlength{\parskip}{5pt}
\setlength{\emergencystretch}{3em}
\begin{document}
\section*{List chromatic number and maximum degree of bipartite graphs}\label{3245}
{\small UnsolvedMath ID 3245 · OPG-47358 · Graph Theory · OpenGarden}
\subsection{Definitions and mathematical statement}
Let $G=(V,E)$ be a finite simple bipartite graph: there is a partition
$V=A\sqcup B$ such that every edge has one end in each part. Let
$\Delta(G)$ be the largest vertex degree. A graph is $k$-choosable if, for
every assignment of a $k$-element set $L(v)$ of colours to each vertex,
there is a proper colouring $c$ with $c(v)\in L(v)$. The least such $k$
is the list chromatic number $\chi_\ell(G)$; for a nonempty edgeless graph
it equals one.

The logarithmic-bound conjecture asks whether there exists an absolute
constant $C>0$ such that, for every finite bipartite graph with
$\Delta(G)\ge2$,
\[
                      \chi_\ell(G)\le C\log\Delta(G).
\]
Here $\log$ denotes the natural logarithm. Changing its base merely
changes $C$. The restriction $\Delta\ge2$ makes the formula meaningful:
$\log1=0$ cannot bound the choice number of an edge. An equivalent
all-degree formulation is that some constant $C'>0$ satisfies
$\chi_\ell(G)\le C'\log(\Delta(G)+2)$ for every nonempty bipartite graph.

The same constant must work independently of the order, the relative
sizes of the two parts, and all overlaps among lists. Bipartite graphs
have ordinary chromatic number at most two, but that does not give a
two-colour list assignment compatible with arbitrary vertex lists.
The target concerns this additional list-colouring difficulty rather
than ordinary two-colourability.
\subsection{Short English statement}
Is the list chromatic number of every bipartite graph bounded by a constant times the logarithm of its maximum degree?
\subsection{Sources}
\begin{itemize}
\item[S1] ULAM.AI and the UnsolvedMath Contributors, supplied problems.json, record 3245 (OPG-47358), OpenGarden collection. Catalogue identity and source transcription; CC BY 4.0.
\url{https://huggingface.co/datasets/ulamai/UnsolvedMath}.
\item[S2] Open Problem Garden, List chromatic number and maximum degree of bipartite graphs. Original problem and discussion; the mathematical formulation below is expanded from the supplied source record.
\url{http://www.openproblemgarden.org/op/list_chromatic_number_and_maximum_degree_of_bipartite_graphs}.
\end{itemize}
\end{document}
